\section{Related Work}
Content moderation refers to the process of detecting user (e.g.,~\cite{gorwa+al:20, kumar+al:24, kolla+al:24}) and AI (e.g.,~\cite{zellers+al:19, lloyd+al:23, fisher+al:24}) generated content that violates laws or policies (cf.~\cite{gorwa+al:20, grimmelmann:15}). The rapid growth of user-generated content, coupled with the vast and evolving range of moderation topics across tasks and use cases, renders manual review impractical, particularly in production environments. This situation drives the increased adoption of unified LLM-based systems for scalable moderation~\cite{kumar+al:24, kolla+al:24, huang:25, levi+al:25}. Recent surveys highlight that uncertainty quantification is a central challenge for deploying such systems reliably at scale~\cite{shorinwa2025surveyuncertaintyquantificationlarge}, and new methods emphasize efficient estimation techniques suitable for production use~\cite{efficient-and-effective}.


The use of LLMs in production moderation systems critically depends on trustworthy confidence estimates, as well-calibrated confidence scores enable safe automation (e.g., selective deferral to human reviewers~\cite{jung+al:24, mozannar2021consistentestimatorslearningdefer}). There is a large body of work on uncertainty quantification in LLMs~\cite{tian+al:23, kadavath+al:22, lin+al:22, kuhn+:23, xiong+al:23, jiang+al:21, manakul+al:23, zhang+al:22, shorinwa2025surveyuncertaintyquantificationlarge,barshalom2025tokenprobabilitieslearnablefast, watson2023explaining, kendall+al:17, farr+al:25, kapoor2024large, cohen2024i, efficient-and-effective, liu+al:24}. Prior work has examined verbalized confidence~\cite{tian+al:23, kadavath+al:22, lin+al:22, xiong+al:23}, but empirical studies consistently find such methods poorly calibrated and highly sensitive to prompting and model type~\cite{kuhn+:23, xiong+al:23}. In addition, multi-sample methods are computationally expensive at production scale. These findings demonstrate that a single, universal uncertainty metric is unreliable and motivate a more robust, multifaceted approach for model failure prediction. To address this, alternative strategies have been explored, including logit-based calibration~\cite{jiang+al:21}, reflection or self-consistency signals~\cite{manakul+al:23}, and reinforcement prompt learning~\cite{zhang+al:22}. Some of them, particularly logit-based information, are used in this work. 

In this work, we propose a novel framework that leverages a diverse set of predictors to improve both calibration and cost-aware decision-making in moderation. Our feature set fuses logit-based features shown to be effective for textual calibration~\cite{jiang+al:21} with additional predictors—including multimodal and uncertainty-attribution signals—that extend applicability beyond text-only settings. Our framework is evaluated across diverse closed- and open-source LLMs on two complementary moderation datasets, including a multimodal benchmark~\cite{levi+al:25}, and employs a learned classifier to produce calibrated confidence estimates.

Another line of work closely related to ours is query performance predictors (QPP) for ad hoc document retrieval~\cite{carmel+Elad:10}. These methods were devised to predict the effectiveness of search results in the absence of relevance judgments. It is common to divide them into two groups: (i) pre-retrieval predictors, which operate prior to retrieval time and often utilize corpus statistics~\cite{hauff+al:08}, and (ii) post-retrieval predictors, which leverage information from the retrieved list~\cite{shtok+al:16,shtok+al:12,tao+al:14,zhou+croft:07}. 
%Supervised prediction methods were also used \cite{arabzadeh+al:21}. 
Recently, and with the advance of LLMs, a new category has emerged: post-generation predictors, which utilize information extracted after LLM generation (e.g., next-token distribution~\cite{huly+al:25}). Pre-retrieval, post-retrieval, and post-generation predictors have also been used to estimate the effectiveness of RAG-based LLM systems~\cite{huly+al:25}. Motivated by this line of work, we incorporate both post-retrieval and post-generation predictors into our feature set, extending them with moderation-oriented attribution features for cost-aware escalation.





